%%%%%%%%%%%%%%%%%%%%%%%%%%%%%%%%%%%%%%%%%%%%%%%%%%%%%%%%%%%%%%%%%%%%%%
% Lớp 9 | Chương IV. Hệ thức lượng trong tam giác vuông | Tỉ số lượng giác của góc nhọn
% Dạng 9H4?11-3: Tìm góc khi biết tỉ số lượng giác, so sánh
%%%%%%%%%%%%%%%%%%%%%%%%%%%%%%%%%%%%%%%%%%%%%%%%%%%%%%%%%%%%%%%%%%%%%%


% nguồn: dan_tu_AI (nhập 17/11/2026)
\begin{ex}%[9H4N11-3]%[ma:1cb407be]
	Cho góc nhọn $\alpha$ thỏa mãn $\cos\alpha=\dfrac12$. Số đo góc $\alpha$ là
	\choice
	{\True $60^\circ$}
	{$30^\circ$}
	{$45^\circ$}
	{$90^\circ$}
	\loigiai{
		Ta có $\cos 60^\circ=\dfrac12$ nên $\alpha=60^\circ$.
	}
\end{ex}

% nguồn: dan_tu_AI (nhập 17/11/2026)
\begin{ex}%[9H4H11-3]%[ma:018e4ac2]
	Sắp xếp các tỉ số lượng giác $\sin 20^\circ$, $\cos 50^\circ$, $\sin 35^\circ$ theo thứ tự tăng dần, ta được
	\choice
	{$\cos 50^\circ<\sin 35^\circ<\sin 20^\circ$}
	{$\sin 20^\circ<\cos 50^\circ<\sin 35^\circ$}
	{$\sin 35^\circ<\sin 20^\circ<\cos 50^\circ$}
	{\True $\sin 20^\circ<\sin 35^\circ<\cos 50^\circ$}
	\loigiai{
		$\cos 50^\circ=\sin 40^\circ$ (hai góc phụ nhau).\\
		Khi góc nhọn tăng thì sin tăng, mà $20^\circ<35^\circ<40^\circ$ nên $\sin 20^\circ<\sin 35^\circ<\sin 40^\circ=\cos 50^\circ$.
	}
\end{ex}

% nguồn: dan_tu_AI (nhập 17/11/2026)
\begin{ex}%[9H4V11-3]%[ma:c46cf740]
	Một đường dốc dành cho xe lăn dài $6\text{ m}$ (tính theo mặt dốc), điểm cao nhất của dốc cao hơn mặt đất $1\text{ m}$. Góc tạo bởi mặt dốc và mặt đất xấp xỉ bằng (làm tròn đến độ)
	\choice
	{$9^\circ$}
	{$80^\circ$}
	{\True $10^\circ$}
	{$6^\circ$}
	\loigiai{
		Mặt dốc là cạnh huyền ($6\text{ m}$), chiều cao $1\text{ m}$ là cạnh đối của góc nghiêng $\alpha$.\\
		$\sin\alpha=\dfrac16\Rightarrow\alpha\approx 9{,}6^\circ\approx 10^\circ$ (dùng máy tính cầm tay).
	}
\end{ex}

% nguồn: dan_tu_AI (nhập 17/11/2026)
\begin{ex}%[9H4V11-3]%[ma:e1d3146d]
	Một cột cờ thẳng đứng cao $6\text{ m}$. Vào một thời điểm, bóng của cột cờ trên mặt đất nằm ngang dài $6\sqrt3\text{ m}$. Gọi $\alpha$ là góc tạo bởi tia nắng mặt trời với mặt đất. Các khẳng định sau đúng hay sai?
	\choiceTF
	{\True $\tan\alpha=\dfrac{1}{\sqrt3}$}
	{\True $\alpha=30^\circ$}
	{\True Khi góc tạo bởi tia nắng với mặt đất là $45^\circ$ thì bóng cột cờ dài $6\text{ m}$}
	{Khi góc tạo bởi tia nắng với mặt đất tăng từ $30^\circ$ lên $45^\circ$ thì bóng cột cờ ngắn đi hơn $5\text{ m}$}
	\loigiai{
		\begin{itemchoice}
			\itemch Đúng. $\tan\alpha=\dfrac{\text{chiều cao cột}}{\text{độ dài bóng}}=\dfrac{6}{6\sqrt3}=\dfrac{1}{\sqrt3}$.
			\itemch Đúng. $\tan 30^\circ=\dfrac{1}{\sqrt3}$ nên $\alpha=30^\circ$.
			\itemch Đúng. Gọi $\ell$ là độ dài bóng: $\tan 45^\circ=\dfrac{6}{\ell}\Rightarrow\ell=6\text{ (m)}$.
			\itemch Sai. Bóng ngắn đi $6\sqrt3-6\approx 4{,}39\text{ (m)}<5\text{ m}$.
		\end{itemchoice}
	}
\end{ex}

% nguồn: dan_tu_AI (nhập 17/11/2026)
\begin{bt}%[9H4V11-3]%[ma:4fdc09d5]
	Một bạn thả diều, dây diều được kéo căng thẳng và dài $50\text{ m}$. Đầu dây ở tay bạn đó, con diều cao hơn tay $30\text{ m}$. Tính góc tạo bởi dây diều và phương nằm ngang (đơn vị: độ, làm tròn đến hàng đơn vị).
	\shortans{37}
	\loigiai{
		Dây diều là cạnh huyền ($50\text{ m}$), độ cao $30\text{ m}$ là cạnh đối của góc $\alpha$ cần tìm.\\
		$\sin\alpha=\dfrac{30}{50}=0{,}6\Rightarrow\alpha\approx 36{,}87^\circ\approx 37^\circ$.
	}
\end{bt}

% nguồn: dan_tu_AI (nhập 17/11/2026)
\begin{bt}%[9H4C11-3]%[ma:e79a692c]
	Mặt cắt của một mái nhà hình chữ A là tam giác cân $ABC$ với đáy $BC=8\text{ m}$ và chiều cao từ đỉnh $A$ đến $BC$ là $3\text{ m}$. Tính số đo góc $BAC$ ở đỉnh mái (đơn vị: độ, làm tròn đến hàng đơn vị).
	\shortans{106}
	\loigiai{
		Gọi $H$ là chân đường cao kẻ từ $A$. Tam giác $ABC$ cân tại $A$ nên $H$ là trung điểm $BC$, $BH=4\text{ m}$.\\
		Tam giác $ABH$ vuông tại $H$: $\tan B=\dfrac{AH}{BH}=\dfrac34\Rightarrow\widehat{B}\approx 36{,}87^\circ$.\\
		$\widehat{C}=\widehat{B}$ nên $\widehat{BAC}=180^\circ-2\widehat{B}\approx 180^\circ-73{,}74^\circ=106{,}26^\circ\approx 106^\circ$.
	}
\end{bt}
